\documentclass[10pt]{article}
\usepackage[a4paper,margin=25mm]{geometry}
\usepackage{amsmath,amssymb,amsthm,mathrsfs,longtable,hyperref}
\setlength{\emergencystretch}{3em}
\newtheorem{theorem}{Theorem}[section]
\newtheorem{lemma}[theorem]{Lemma}
\newtheorem{proposition}[theorem]{Proposition}
\newtheorem{corollary}[theorem]{Corollary}
\newcommand{\PSH}{\operatorname{PSH}}
\newcommand{\FS}{\operatorname{FS}}
\newcommand{\MA}{\operatorname{MA}}
\newcommand{\vol}{\operatorname{vol}}
\newcommand{\an}{\mathrm{an}}
\newcommand{\EE}{\mathcal E^1_\theta}
\title{Single-net regularization and exact negligible-set transfer in big classes\\
\large An incomplete standalone paper for the unchanged original P16.5}
\author{Research draft relative to the two precise caller grants}
\date{6 October 2026}
\begin{document}
\maketitle
\begin{abstract}
We reconstruct genuine decreasing closure, envelopes, and normalized
compactness on the original carrier from the actual ample theorem and the
single resolved multiplier-ideal net granted with the assignment. The
flattening argument keeps an explicit weakened pole floor and a separate
positive guard. We then prove that every raw-supremum discrepancy of the
retained cone is covered by a genuine big-class psh pole. The raw supremum
and its usc regularization are identified literally at every point in a
common actual ample embedding; the pole returns by an actual section
tensor identity. Measure nullness follows only under explicitly stated
finite-energy and contraction inputs. A separate finite polarization
calculation identifies exactly what bounded convex-combination convergence
would supply for mixed products, and an algebraic example excludes replacing
that input by ordinary volume interpolation. We have not constructed every
unbounded mixed or weighted product, or established the complete equation.
The canonical result is candidate=false.
\end{abstract}

\section{Exact task, domain, and authority}
The original P16.5 asks for the finite-energy non-Archimedean equation
in big or pseudoeffective nonample classes, with the psh cone, reference
singularity, energy, and total mass specified, and with a replacement for
probability normalization at zero volume. The retained first domain is
every smooth integral projective $k$-variety $X$, of every dimension
$n\geq1$, over every algebraically closed characteristic-zero field
with the trivial absolute value, and every actual big rational line bundle
$L\in\operatorname{Pic}(X)\otimes\mathbb Q$. Write
$K=X^\an$, $\theta=c_1(L)$, and $V=\vol_X(L)>0$. No countability,
finite-generation, invariance, skeleton, or support hypothesis is added.
All supported semivaluations stay in the compact Hausdorff space $K$.
The scalar reference of a bundle is its canonical trivial metric.

Only G0 and G1 are caller mathematical grants. G0 is the supplied
complete actual ample rational-bundle theorem with its exact retained
field and geometric hypotheses. G1 is the supplied resolved relative
generation theorem and its one generically finite FS net. We use G0 on
the original smooth $X$ in actual ample classes. We use G1 on smooth
equivariant modifications with the stated actual vertical Cartier data,
full relative systems, and actual ample twist. No extension of either
grant to big-class energy or equations is presumed. Other manuscripts
are unreviewed mathematical data; their propositions are not premises
of the positive results below.

Ordinary algebraic background is Kodaira's ample-plus-effective
decomposition of a big rational bundle, openness and numerical invariance
of ampleness, eventual global generation, and characteristic-zero
equivariant model resolution. Section 5 uses a separately retrieved
primary theorem in its actual ample scope. We distinguish that application
from both grants. We do not claim a recursive new proof of algebraic
resolution or the whole ample capacity theory.

An actual FS scalar on a rational bundle $B$ is
\[
 h=\frac1m\max_j\bigl(\log\|s_j\|_{\rm triv}+\lambda_j\bigr),
 \quad 0\ne s_j\in H^0(X,mB),\quad\lambda_j\in\mathbb R.
 \tag{1}
\]
Common zeros are permitted. These are continuous into
$\mathbb R\cup\{-\infty\}$ and have finite value at the generic
trivial valuation $v_0$. A genuine $\theta$-psh scalar is the all-point
decreasing limit of \emph{one} such net on actual $B_i$ with
$c_1(B_i)\to\theta$, and with finite generic endpoint. This is our
literal cone $\PSH(\theta)$, rather than a twice-closed family.
Every bundle exponent is cleared before taking sections.

The intended energy and target data, not constructed here in their
whole required scope, are the section-volume outer energy
\[
 E(u)=\inf_{f\in C(K),\ f\geq u}\mathscr E(f),\qquad
 \EE=\{u\in\PSH(\theta):E(u)>-\infty\},\qquad R=P_\theta(0),
 \tag{2}
\]
and positive Radon targets of raw mass $V$ such that
\[
 \nu\{R=-\infty\}=0,\qquad
 D_\nu=\sup_{u\in\EE,\ u(v_0)=0}
 \left(E(u)+\int(R-u)\,d\nu\right)<\infty.                 \tag{3}
\]
Differences in (3) are taken where $R$ is finite; no absolute
integrability of $R$ is silently required. The genuine full space in
(2), rather than only its bounded relative part, is an outstanding
operator obligation. The reference $R$ itself is constructed below.

\section{Exact scalar operations and bounded section density}
At the valuation-ring center of each original analytic point, a
generating section of an actual globally generated bundle is a unit.
All canonical section norms are at most one. Consequently the
weight-zero maximum of a generating system is exactly zero at every
point, including supported points. Tensoring with this system changes
the actual bundle but preserves the scalar everywhere.

\begin{lemma}[Common actual ample embedding]\label{embed}
Choose an actual ample rational $\Omega$ such that $\Omega-L$ is
ample. Every $u\in\PSH(\theta)$ is the same all-point scalar function
in the genuine ample cone of $\Omega$.
\end{lemma}
\begin{proof}
For a tail of a defining net, $\Omega-B_i$ is ample. Clear powers
and tensor with its generating system to obtain the same scalar on
$\Omega$. To remove common zeros, replace the scalar stage by
$\max(h_i,-N)$ using generators of a power of $\Omega$ with weight
$-mN$. The product-directed net $(i,N)$ is decreasing, has fixed
actual bundle $\Omega$, and has all-point infimum $u$. Its generic
endpoint remains finite. It is an actual ample defining net within G0.
No infinite value was subtracted and no modification of $X$ was used.
\end{proof}

In particular the generic value equals the supremum; all divisorial
values are finite and determine all-point order by G0. Normalized
functions have finite uniform coordinate bounds in the one compact
normalized ample cone. The same embedding works at any fixed real
numerical class $\eta$ using an actual ample $\Omega$ with
$c_1(\Omega)-\eta$ ample, without requiring a rational representative
of $\eta$.

Finite maxima in the same class have single defining nets: put two
incoming stages on the actual $L+\varepsilon A$ whenever both actual
differences are ample, and unite their section systems. Valid triples
of indices and rational $\varepsilon>0$ are directed by choosing a
smaller $\varepsilon$ and sufficiently late indices. The scalars
decrease and their classes tend to $\theta$.
For a real convex coefficient $0<r<1$, put incoming scalars below a
common finite $C$, choose positive rational $r_j\uparrow r$ and
$s_j\uparrow1-r$, and use
\[
 r_j(h_i-C)+s_j(g_l-C)+C.
\]
This is a decreasing product net of actual FS scalars on
$r_jB_i+s_jB'_l$. Its class tends to $\theta$ and its generic endpoint
is finite. Constants, rational scaling, and integer tensoring are direct.

Fix one actual Kodaira identity $L=aA+E_0$, with $a>0$ rational,
$A$ ample, and $E_0$ effective rational Cartier. Set
$W=\log\|s_{E_0}\|\leq0$, $Z=E_0^\an$, and $U=K\setminus Z$.
Tensoring its section with a generating power of $aA$ realizes $W$
on the actual $L$. It is finite continuous on $U$.

\begin{lemma}[Bounded minorants]\label{density}
For every bounded continuous $f$ and every genuine $u\leq f$,
\[
 u\leq Q_L(f):=\sup\{h\in\FS(L):h\leq f\}\quad\text{on }U.
 \tag{4}
\]
This statement does not presume decreasing closure or a genuine envelope.
\end{lemma}
\begin{proof}
Take its literal net $h_i\downarrow u$ on $B_i$. Fix rational
$0<t<1$ and $\epsilon>0$ before taking a tail. The actual bundle
$L-(1-t)B_i-t(L-aA)$ is eventually ample. The nested compact sets
$\{h_i\geq f+\epsilon\}$ have empty intersection, so a whole tail
has $h_i<f+\epsilon$ on all of $K$. Set $C=\|f\|_\infty$.
Tensor the sections with the zero-weight system of the actual ample
difference. This gives the actual $L$-FS minorant
\[
 (1-t)(h_i-\epsilon)+t(W-C)\leq f
\]
at every original point. At $x\in U$ with $u(x)$ finite, take the
incoming endpoint, then $t\downarrow0$, then $\epsilon\downarrow0$.
The value of $W(x)$ is finite, giving (4). If $u(x)=-\infty$, the
inequality is immediate. Supported zeros satisfy the minorant inequality
literally as extended values.
\end{proof}

\begin{proposition}[First genuine model envelope]\label{model}
For every rational continuous model obstacle $f$, its all-point raw
envelope $P_\theta(f)$ is genuine. In particular $R=P_\theta(0)$
is genuine, $R(v_0)=0$, and normalized genuine $u$ satisfy $u\leq R\leq0$.
\end{proposition}
\begin{proof}
Represent $f$ on a resolved equivariant model in G1's exact signature.
That grant supplies one genuine endpoint $\Psi$ with
$Q_L(f)\leq\Psi\leq f$. Lemma~\ref{density} puts every genuine
competitor below $\Psi$ off $Z$, hence at every divisorial valuation.
Embed the two already genuine functions into the same actual ample
$\Omega$. G0 order determination gives the inequality at every point
of $K$. The endpoint $\Psi$ is itself a competitor, so it is exactly
$P_\theta(f)$, without taking any second decreasing closure.
For $f=0$, $W$ is a genuine minorant with generic value zero, proving
the reference normalization and the stated maximality.
\end{proof}

For any actual big rational $B$, the same result constructs $R_B$.
A late G1 stage for its zero obstacle, on the actual $S=B+C/m$,
can be normalized to a scalar $\psi_B\leq0$. Its original generic
shift $c_B$ is finite and
\[
 R_B\leq\psi_B+c_B.                                      \tag{5}
\]
We do not assert that the separately normalized stages form a
decreasing net; only each fixed reference in (5) is used next.

\section{Single-net decreasing closure across old poles}
\begin{lemma}[Compact pole comparison]\label{polecompare}
Let $\psi\leq0$ be an actual FS scalar and let an actual FS scalar
$F$ satisfy $F\geq(1-\delta)\psi-k$, where $0<\delta<1$.
For usc $h\leq\psi+C$, the comparison $h-F$, assigned value
$-\infty$ at poles of $F$, is usc. Decreasing such $h$ give decreasing
comparison functions. If their infimum is negative everywhere, a whole
tail of the comparisons is negative everywhere.
\end{lemma}
\begin{proof}
Where $F$ is finite the comparison is usc and at most
$\delta\psi+C+k$. A pole of $F$ forces $\psi=h=-\infty$.
Extended continuity of $\psi$ makes this upper bound tend to
$-\infty$ near that point, proving usc of the assigned comparison.
This is an extension of a comparison, not arithmetic
$(-\infty)-(-\infty)$. The nonnegative superlevels are nested compact
sets. Empty intersection and directedness make one such set, then
all later ones, empty.
\end{proof}

\begin{theorem}[Genuine decreasing closure]\label{closure}
If $u_i\in\PSH(\theta)$ decrease all-point to $u$ with $u(v_0)$
finite, then $u$ has one all-point decreasing FS net with classes
tending to $\theta$.
\end{theorem}
\begin{proof}
Restrict to a tail $u_i\leq M$. For rational $\eta>0$ let
$T_\eta=L+\eta A$. Choose $0<\epsilon<\eta/4$ and
$B=L+\epsilon A$. Use (5) to choose a normalized reference stage
$\psi\leq0$ on an actual $S$ sufficiently close to $B$ that
$T_\eta-S$ is ample. Choose small rational $0<\delta<1$ with
$T_\eta-(1-\delta)S$ still ample.

A special stage is an actual FS scalar on the specified $T_\eta$
with complete reference data as above and with
\[
 F\geq u+\kappa,\qquad F\geq(1-\delta)\psi-k,
 \tag{6}
\]
where $\kappa>0$ is rational and $k$ finite. Such a stage exists:
tensor a late defining majorant of one $u_i$, shifted by $\kappa$,
and the pole floor into $T_\eta$, then take their maximum.

For every fixed old $B_l,\psi_l$, all $u_i$ embed as scalars in
$\PSH(c_1(B_l))$: their incoming stages eventually have actual ample
difference to $B_l$. Late stages $h$ of a fixed $u_i$ do likewise.
Their generic values are eventually at most $M+1$. The model
reference maximality, valid before closure, gives
\[
 h\leq R_{B_l}+M+1\leq\psi_l+c_l+M+1.                    \tag{7}
\]
The same bound holds for all $u_i$. This is the common old-pole
control needed for Lemma~\ref{polecompare}.

Given finitely many old special stages $F_l$ with guards $\kappa_l$,
choose new $\eta$ arbitrarily small below all their $\epsilon_l$.
Choose its new reference stage $\psi$ on $S$ sufficiently close to
$L+\epsilon A$ that $B_l-S$ is ample for every $l$. Thus
$\psi\leq R_{B_l}\leq\psi_l+c_l$. Choose
$0<\delta<\min_l\delta_l$ satisfying the new target ample gap.
Where $F_l$ is finite, the new floor $b=(1-\delta)\psi-k$ satisfies
\[
 b-F_l\leq(\delta_l-\delta)\psi_l+(1-\delta)c_l+k_l-k.
\]
Since $\psi_l\leq0$, one large $k$ gives $b\leq F_l$ for all $l$.
At an old pole, the reference inequality forces $b=-\infty$, so
this inequality is literal there too.

Choose $0<\kappa<\min_l\kappa_l/4$. Apply
Lemma~\ref{polecompare} to the decreasing $u_i+2\kappa$ using (7).
Its limit comparison is negative: at a finite old value it is at
most $2\kappa-\kappa_l<0$, and at an old pole it is assigned
$-\infty$. One common late $i$ satisfies every old comparison.
Apply the lemma again to the literal defining tail of this fixed
$u_i$, shifted by $\kappa$, with (7). Choose one stage $h$, late
enough to fit the new actual $T_\eta$, such that
$h+\kappa\leq F_l$ for all $l$. Then
\[
 F=\max(h+\kappa,b)\leq F_l\quad(\text{all }l),
 \qquad F\geq u+\kappa.
\]
It is a new special stage. Order all special stages by scalar
decrease and by $\eta$ decrease. This is a directed net whose class
tends to $\theta$.

Its infimum is $u$ at every point. Given $x$ and finite $t>u(x)$,
choose $u_i(x)<t$ with margin, choose a target gap $\eta$ before
a sufficiently late defining stage $h$ fitting it, and choose its
positive guard sufficiently small that $h(x)+\kappa<t$. A large
$k$ makes the floor below $t$ at $x$, or it already has value
$-\infty$ there. This constructs a special stage below $t$ at $x$;
directedness can impose every old constraint afterwards. The argument
works for every finite $t$ when $u(x)=-\infty$. All special stages
stay above $u$ and the stipulated generic endpoint is finite.
\end{proof}

\section{Genuine envelopes and original compactness}
\begin{theorem}[Envelopes and raw-supremum regularization]\label{envelopes}
Every upper bounded usc obstacle $f$ with a genuine generically finite
minorant has a genuine all-point envelope $P_\theta(f)$.
For any nonempty uniformly upper bounded family $(u_\alpha)$ in
$\PSH(\theta)$, the literal function
$(\sup_\alpha u_\alpha)^*$ is genuine $\theta$-psh.
\end{theorem}
\begin{proof}
On a compact Hausdorff space an upper bounded usc function, including
one with poles, is the pointwise infimum of continuous majorants.
For each point and strict finite upper value, use a continuous cutoff
inside the open strict sublevel neighborhood, and a global upper
bound outside that neighborhood. Uniform rational model density from
G0 and finite minima give a directed decreasing net of rational
model majorants with the same infimum. Their envelopes are genuine
by Proposition~\ref{model}, decrease, and dominate the stipulated
minorant at the generic point. Theorem~\ref{closure} makes their
endpoint genuine. It is below $f$ at every point; every original
competitor is below each envelope. Thus the endpoint is exactly
the raw $P_\theta(f)$.

Let $S=\sup_\alpha u_\alpha$ and $f=S^*$. There is a finite generic
lower bound from any one family member and a finite global upper bound.
Each member is a minorant of $f$, so $P_\theta(f)$ dominates $S$.
It is usc, hence dominates the least usc majorant $S^*$. The obstacle
inequality gives the opposite inequality. Therefore
$P_\theta(S^*)=S^*$ literally at every original point, not merely
on divisorial coordinates.
\end{proof}

For continuous $f$, Lemma~\ref{density} now proves
$P_\theta(f)=Q_L(f)$ on $U$. The left side is usc, the right side
is lsc, and both are finite there, so the envelope is continuous on
this ordinary open set. This equality is not asserted on $Z$.

\begin{theorem}[Normalized compactness]\label{compact}
The normalized genuine cone $\{u\in\PSH(\theta):u(v_0)=0\}$ is compact
in its divisorial topology.
\end{theorem}
\begin{proof}
Embed it into the compact normalized actual ample cone of $\Omega$.
Consider a net in the big cone converging in that ample cone to $a$.
Let $S_i^*=(\sup_{j\geq i}u_j)^*$. By Theorem~\ref{envelopes} these
are genuine $\theta$-psh, decrease, and have generic value zero.
Theorem~\ref{closure} gives their all-point infimum $b$ in the
normalized big cone. The primary finite separator theorem in the
actual ample cone, BJ Theorem 5.6 and Lemma 5.7, gives
$S_i^*(z)=\sup_{j\geq i}u_j(z)$ at each divisorial $z$. Hence
$b(z)=a(z)$ at all those coordinates. G0 all-point determination
in $\Omega$ gives $b=a$ everywhere. The normalized big cone is
therefore a closed subset of the compact ample cone.

The external separator result is applied only after the literal
common ample embedding. Its pertinent primary statement and full
proof have been independently read and saved. No big energy compactness
or equation is used in this proof.
\end{proof}

If the topology is indexed only by rational divisorial scalings, the
conclusion is unchanged. Along a ray $t\operatorname{ord}_E$, an FS
stage is a maximum of affine functions with nonpositive slopes. Its
finite decreasing endpoint is convex and decreasing on $t>0$.
Uniform ample coordinate bounds on a larger compact interval give
a uniform Lipschitz bound on each smaller interval, using the usual
secant inequalities for convex functions. Values on rational scalings
therefore determine both real-scaled values and their convergence.
This adds no field-countability assumption.

\section{Literal negligible-set transfer and genuine big poles}
For a bounded-above family in one single-net numerical cone, put
$S=\sup_\alpha u_\alpha$ pointwise and define its discrepancy
$D=\{S<S^*\}$. A negligible subset means any subset of such a
discrepancy, with the same original analytic topology.

\begin{theorem}[All-point negligible-to-big-pole transfer]\label{neglig}
Every such negligible subset $D'\subset K$, including discrepancies
from any fixed real numerical class with the single-net convention,
is contained in $\{q=-\infty\}$ for a genuine
$q\in\PSH(\theta)$ on the original $K$.
In particular this applies to every discrepancy of the retained big
cone. The witness has one actual section defining net.
\end{theorem}
\begin{proof}
Choose one actual ample $\Omega$ whose numerical difference from
the input family's class is ample. Lemma~\ref{embed} (with that
class) embeds each member as the \emph{identical scalar function} at
every point of $K$. Thus the raw supremum is identical, and its usc
regularization is identical, because the underlying topology and all
pointwise values are identical. Explicitly its value is the infimum
over neighborhoods $O$ of $x$ of $\sup_{y\in O}S(y)$ in both
descriptions. There is no equality inferred from divisorial density.
Hence $D'$ is an actual ample negligible subset in the exact primary
Definition 5.1 convention.

Apply BJ v4 Theorem 13.17 in that ample class: $X$ is irreducible,
$\Omega$ is ample, and EP is granted by G0. There is an ample-psh
$p$ with $p=-\infty$ on $D'$. Normalize $p(v_0)=0$.
This source theorem is an external premise in its ample domain,
not an assertion that negligible and polar mean the same by definition.
The whole pertinent statement/proof, including Lemma 13.19, has
been independently read; its capacity/Choquet dependencies are
documented in the evidence appendix.

Choose a positive rational $t$ so small that $L-t\Omega$ is big.
A nonzero section of a cleared positive multiple supplies an actual
effective rational divisor $E$ and an actual rational bundle identity
$L=t\Omega+E$. Let $q_E=\log\|s_E\|\leq0$.
G0 supplies a literal decreasing fixed-bundle FS net $p_i\downarrow p$
on $\Omega$. Clear $t$, the incoming section powers, and the divisor
power. Tensor each section of $p_i$ with the canonical section of
$E$ to obtain an actual $L$-FS stage with scalar
\[
 q_i=t p_i+q_E.
\]
This is one decreasing net. At $q_E$ poles it has value
$-\infty$ at every stage, and elsewhere its infimum is
$tp+q_E$. Thus its all-point endpoint $q$ is genuine in
$\PSH(\theta)$, is finite at $v_0$, and is $-\infty$ on $D'$.
This is an actual tensor identity on the original $X$, with no
arbitrary lift and no subtraction of supported poles.
\end{proof}

Conversely a genuine big-class pole is an ample pole by
Lemma~\ref{embed}. Thus the big and ample pole-covering collections
agree for this retained big domain. This assertion alone gives no
finite-energy membership for the transferred $q$: its chosen fixed
effective divisor can exceed the minimal singularity type. In
particular it does not give an unconditional output-nullness theorem.

Theorem~\ref{envelopes} supplies genuine big regularization independently
of this source application. For a big input family, $S^*$ is the same
all-point regularized scalar in both cones and is genuinely in the
big cone. We have proved precisely the all-point statement missing
from a merely divisorial identification.

\section{The remaining finite-energy nullness antecedents}
The following lemma makes the logical boundary of the new negligible
argument explicit. In this section only, assume that the energy (2)
has been independently established on the entire genuine cone, with
concavity, monotonicity, decreasing continuity, $E(R)=0$ and raw
constant law $E(u+c)=E(u)+cV$. Assume also its genuine contraction
\[
 B((1-r)R+ru)\leq C r^2 B(u),\quad B(u)=-E(u),
 \quad u(v_0)=0,\quad0<r<1,                              \tag{8}
\]
where $C$ is independent of $u$. These are additional unproved
energy inputs of this paper, not consequences of Section 5 or new grants.

\begin{lemma}[One genuine E1 pole covers each psh pole]\label{e1pole}
Under these energy inputs, for every normalized
$q\in\PSH(\theta)$ there is $p\in\EE$ such that
\[
 \{q=-\infty,\ R> -\infty\}\subset\{p=-\infty\}.
\]
\end{lemma}
\begin{proof}
Let $N_j=16^j$, $r_j=8^{-j}$, and $\lambda_j=2^{-j}$.
The genuine truncation $q_j=\max(q,R-N_j)$ is normalized and
$R-N_j\leq q_j\leq R$. Monotonicity and the raw constant law
give $B(q_j)\leq VN_j$, so $q_j\in\EE$.
Set $a_j=(1-r_j)R+r_jq_j$. Then (8) gives
$B(a_j)\leq CV4^{-j}$. The actual convex functions
\[
 p_m=\sum_{j\leq m}\lambda_j a_j+2^{-m}R
\]
decrease, have generic value zero, and have energy at least
$-CV\sum_{j\geq1}8^{-j}=-CV/7$ by concavity.
Theorem~\ref{closure} gives a genuine all-point endpoint; the
assumed decreasing energy continuity puts it in $\EE$.
Where $q=-\infty$ and $R$ is finite, $R-q_j=N_j$ and the exact
finite relative sum is
$R-p_m=\sum_{j\leq m}\lambda_jr_jN_j=m$. Thus $p=-\infty$ there.
\end{proof}

\begin{corollary}[Conditional full negligible nullness]\label{null}
Under these energy inputs, any finite positive Radon measure missing
$\{R=-\infty\}$ and integrating every genuine relative E1 test
has outer measure zero on every negligible subset of Section 5.
\end{corollary}
\begin{proof}
Theorem~\ref{neglig} supplies $q$. Lemma~\ref{e1pole} supplies one
actual E1 pole covering its pole outside the excluded reference poles.
Its nonnegative relative cost has finite integral, so the measure of
that Borel infinite-cost set is zero. The reference-pole set is also
Borel and null. Their union is a Borel null superset of the original
negligible set, even when the latter is not itself assumed Borel.
\end{proof}

This implication treats nonalgebraic raw-supremum discrepancies.
It cannot be used to establish its own energy or all-test output
integrability antecedents. In particular we have not proved that
the proposed outer-energy outputs satisfy those antecedents.
The finite-dual targets (3) would satisfy the all-test condition
directly once (2) is established, since each genuine test has finite
cost. The operator image inclusion remains a separate assertion.

\section{Finite polarization: an exact implementation and its boundary}
On any one finite common model, with nef slots $T_1,\ldots,T_n$
of the same generic bundle, the raw mixed measure has component
coefficient $b_E(T_1|_E\cdots T_n|_E)$ at $v_E$. These are ordinary
finite intersections. Their positive finite nef interpretation follows
from G0 on the actual ample $A+\epsilon H$ and the finite polynomial
limit as rational $\epsilon\downarrow0$; the scalar models are kept
fixed during that limit. Multilinearity gives
\[
 T_1\cdots T_n=\frac1{n!}
 \sum_{\varnothing\ne S\subset\{1,\ldots,n\}}
 (-1)^{n-|S|}\Bigl(\sum_{i\in S}T_i\Bigr)^n.             \tag{9}
\]
For raw model measures this is equivalently
\[
 \mu(T_1,\ldots,T_n)=\frac1{n!}
 \sum_{S\ne\varnothing}(-1)^{n-|S|}|S|^n\,
 \MA_{\rm raw}\left(|S|^{-1}\sum_{i\in S}T_i\right).
 \tag{10}
\]
Every central multiplicity is retained in each diagonal measure.
There is no division by the generic volume in this formula.

\begin{proposition}[What diagonal convergence would prove]\label{polarization}
Suppose, for a fixed finite list of genuine endpoints and admissible
shared moving approximants, the diagonal raw measures converge for
each of the finitely many rational barycenters in (10), with the
correct actual endpoint operator. Then the finite mixed measures
converge weakly to the expression (10) of those endpoint outputs.
That limit is positive, has the limiting raw mixed mass, and is
independent of any approximant family satisfying the same hypothesis.
\end{proposition}
\begin{proof}
Equation (9) follows by expanding homogeneous monomials of degree
$n$. A monomial whose support is a proper subset cancels in the
alternating superset sum. The only full-support monomial has every
exponent one and coefficient $n!$. This proves (9) for all $n$,
then (10) by homogeneity of the raw diagonal product.
For each fixed continuous test the finite sum (10) passes through
the stipulated diagonal limits. Each finite mixed measure is
positive, so its limiting integral of every nonnegative continuous
test is nonnegative. Riesz positivity gives the positive Radon
limit. The constant test gives its raw mass, and the resulting finite
sum depends only on the endpoint outputs. This proves the claimed
restricted independence. No unbounded test has been passed.
\end{proof}

The exact rational recipe is implemented in \texttt{polarize.py}.
For example in dimension two it is
$2\MA((u_1+u_2)/2)-(\MA(u_1)+\MA(u_2))/2$.
The program gives subset coefficients and checks the polynomial
identity and the raw mass coefficient. The proof is the all-dimension
cancellation argument, not its finite tests.

The hypotheses of Proposition~\ref{polarization} are \emph{not}
proved by the bounded moving-max argument in the supplied ROOT201.
A union of degree-$m$ systems represents a maximum on that same
degree. A rational convex combination generally uses tensor products
on a larger degree; the horizontal moving factor and its full-section
defect must be compared again. Convergence for endpoints and their
maximum alone does not prove the barycenter limits in (10).
The raw positive products of big algebraic classes are not
multilinear substitutes for this step.

Here is a direct algebraic safeguard, with no specialized external
volume-differentiation premise. On $\operatorname{Bl}_p\mathbb P^2$
write $H$ for the pulled hyperplane and $E$ for the exceptional curve.
For rational $-1<a<1$ let $D(a)=H+aE$. In a denominator-cleared
degree, if $a\geq0$, negative intersection with $E$ forces $maE$
as a fixed component; removal gives the whole system $|mH|$.
If $a=-b<0$, sections are degree-$m$ polynomials vanishing to order
$mb$ at $p$. In affine coordinates at $p$ the monomials of total
degree less than $mb$ impose independent conditions, because $mb<m$.
Therefore
\[
 h^0(mD(a))=\begin{cases}
 \binom{m+2}{2},&a\geq0,\\
 \binom{m+2}{2}-\binom{mb+1}{2},&a=-b<0,
 \end{cases}
\quad
 \vol(D(a))=\begin{cases}1,&a\geq0,\\1-a^2,&a<0.\end{cases} \tag{11}
\]
All these bundles are big. The volume function is constant on one
open half-interval and nonconstant on the other, so it is not one
polynomial on this entirely big segment. A quadratic fit at
$a=-1/2,0,1/2$ gives $1+a/4-a^2/2$; at $a=1/4$ it gives
$33/32$ while the true volume is $1$. This rules out obtaining the
required polarization merely by interpolating ordinary algebraic
volume along arbitrary big classes. It does not rule out a valid
same-class finite-energy mixed construction.

\section{Audit of the Borel transfer and remaining original obligations}
We separately checked the fixed-ramp mechanism of the supplied ROOT201.
Its crucial implication is correct under its stated genuine and output
inputs. With canonical $u_i\downarrow u$, $v_i\downarrow v$, a fixed
index $j$, and a fixed ramp vanishing when $u_j-v_j\leq\epsilon$,
an incorrect selection at $i\geq j$ with $u_i\leq v_i$ forces
$u_j-u_i>\epsilon$. Uniform ordered absolute errors therefore bound
the two offending positive measures before a weak limit is taken.
The test is held continuous and fixed through that limit; afterwards
one takes $j$ to its endpoint in absolute L1, then an ordinary
sequence of ramp widths to zero at fixed measures. This is a valid
conditional Borel-selection argument. No openness of the unbounded
winner or weak test by its moving indicator is required.

This check does not discharge the whole inputs. The scalar cone,
all-point envelopes, and compactness are reconstructed in Sections
2--4. We have not independently established the whole raw outer
energy, its unbounded mixed compatibility, contraction, uniform
absolute midpoint estimates, or the output map used by that Borel
argument. In particular finite nef locality needs the genuine finite
model coefficient calculus from G0 before a finite polynomial ample
perturbation; it cannot justify an unbounded nonample passage by itself.

The exact remaining statements for the unchanged original are:
\begin{enumerate}
\item Establish the entire scalar energy (2) and its weak usc,
decreasing continuity, raw mass constant law, genuine contraction,
and every finite moving trace compatibility, retaining lost-volume
errors until the fixed-continuous-test limit.
\item Construct the target-independent original-carrier operator
and prove both continuous energy derivative signs. Identify it
with all bounded rational barycenter approximations needed in (10),
then all mixed products in the full genuine E1 domain.
\item Identify every admissible simultaneous unbounded physical
weighted approximation and its independence. Physical section
functions with horizontal poles are not thereby original E1 tests.
Prove uniform time-deficient tail bounds and caps with every degree,
collar, cost, and raw mass limit in a justified order.
\item Establish all genuine relative E1 output integrability and
contraction. Only then Corollary~\ref{null} supplies full negligible
nullness, rather than just algebraic-divisor nullness.
\item Audit and discharge the hypotheses of the supplied absolute
midpoint, no-jump, attainment, full supporting-plane, and strict
separation arguments. An equation for the outer map alone is
insufficient until the required original operator is identified.
\item Discharge the bounded-max measure input and uniform ordered
errors of the strict Borel argument, then obtain genuine comparison
and uniqueness of every solution modulo constants.
\item Specify a replacement datum for any additional nonzero
zero-volume pseudoeffective range. No such additional range is
established in this paper.
\end{enumerate}
These are limitations of this submission, not a claim of
bibliographical nonexistence. No original P16.5 solution or other
P16-axis solution is asserted.

\section{The additional zero-volume range retained independently}
For the actual numerically trivial rational-bundle range, use the
same varying-actual-bundle cone at class zero, with zero reference.
This result needs only G0 and the literal single-net convention.

\begin{proposition}
In this convention $\PSH(0)=\mathbb R$. The raw zero-class top
operator and raw energy are zero on the entire cone; its finite-energy
space is $\mathbb R$ and its sole positive mass-zero target is the
zero measure. The equation is attained and unique modulo constants.
\end{proposition}
\begin{proof}
For any positive rational $\epsilon$, an actual numerically trivial
$L+\epsilon A$ is ample. Its zero-weight generating scalar is exactly
zero everywhere. These constant scalar stages with classes tending
to zero give one defining net; shifts give all constants.
Conversely normalize $u\in\PSH(0)$ at $v_0$ to zero. Integer
tensoring makes $mu\in\PSH(0)$ for every positive integer $m$.
Embed all of them into one fixed actual ample $\Omega$. Its normalized
compactness and finite continuous divisorial coordinates give
$-C_z\leq mu(z)\leq0$, uniformly in $m$. Thus $u(z)=0$ for
all divisorial $z$, and all-point determination gives $u=0$ on $K$.
The ample perturbations have raw mass $\epsilon^n(A^n)\to0$.
Positivity bounds every continuous test integral by this mass times
its sup norm. Thus the operator limit is zero. Its constant energies
are $c\epsilon^n(A^n)\to0$, giving zero energy on the entire
classified cone. A positive Radon measure with mass zero is literally
zero, and all solutions are constants. Comparison and the two zero
energy derivative signs are immediate. There is no division by zero.
\end{proof}

The nonzero nef class $p_1^*\mathcal O(1)$ on
$\mathbb P^1\times\mathbb P^1$ has zero volume and a nonconstant
bounded generated-system scalar obtained by pulling back the weighted
maximum of the two homogeneous coordinates with weights $0,-1$.
It and zero have vanishing raw top limits: their common positive ample
perturbation masses are $2\epsilon(1+\epsilon)\to0$. Therefore
uniqueness modulo constants fails for that whole zero-volume bounded
cone with the raw top equation. This safeguard is not an additional
nonzero zero-volume equation theorem.

\appendix
\section{Primary reading, preserved inputs, and reproducibility}
The exact input bytes, task object, full family, assignment, and all
18 marked supplied mathematical bodies are preserved in
\texttt{frozen/} and \texttt{supplied/}. Whole-unit inclusive
assignment-LF coordinates and hashes are in \texttt{supplied/index.json}.
Historical paths and hashes inside ledgers are documentary data; no
such path was opened and no historical or private verdict was inspected.

BJ means arXiv:1801.08229v4, PDF dated 24 March 2022. Primary
connector responses were fetched independently in this workspace.
The complete applied finite-separator unit is original renderer
3996--4022. The applied negligible-to-polar unit is 10614--10647,
including its exact ample-EP hypotheses and full pertinent proofs.
Ancillary statement/proof units read are 3863--3898 (polar detection),
10092--10108 (Choquet), 10289--10613 (quasicontinuity and capacity
comparison), and 10738--10847 (outer capacity). The literal
negligible definition is 3921--3928. Countable auxiliary selections
in the ample primary proof are conclusions of that theory, not
countability assumptions on our field, carrier, or given families.

The Addendum is arXiv:2206.07183v2. Its complete bounded-section-density
statement/proof at renderer 128--183 was independently read as source
matching; our proof uses its own actual ample-gap calculation and does
not import the Addendum's general closure. G1 alone supplies its exact
resolved model endpoint. BFJ v2 was opened diagnostically; no new
BFJ theorem is an external premise of this paper. The elementary
volume computation (11) is proved from its explicit section counts.

\texttt{sources/unit-index.json} maps all eight complete selected units
to local inclusive one-based LF coordinates and byte hashes, keeping
the original renderer and PDF page labels separate. No overlap conflict
or missing line remains. The source appendix records complete pertinent
units, not a claim to have read whole primary papers or recursively
all their background proofs. No original downloaded PDF is asserted:
a default-transport request failed with operation not permitted, and
the observation is preserved in \texttt{logs/network-attempt.json}.

Run \texttt{python3 rerun.py} here. It verifies the frozen bytes,
whole-unit slices, primary coordinates, backward note-support order,
and exact finite polarization and volume examples, and compiles this
paper twice with the installed pdfLaTeX. Programs, inputs, commands,
outputs, tool versions, and manifest are retained locally. These
checks are reproducibility and transcription evidence, not mathematical
certification. No resource-heavy calculation was launched; no
operator-provided supervised compute service was exposed. No agents,
credentials, account state, runtime settings, or other campaigns
were read or changed. The current native model identity is not
exposed by these tools and no model switch was attempted.

\begin{thebibliography}{9}
\bibitem{BJ} S. Boucksom and M. Jonsson,
\emph{Global pluripotential theory over a trivially valued field},
\href{https://arxiv.org/pdf/1801.08229v4}{arXiv:1801.08229v4}.
\bibitem{ADD} S. Boucksom and M. Jonsson,
\emph{Addendum to global pluripotential theory over a trivially valued field},
\href{https://arxiv.org/pdf/2206.07183v2}{arXiv:2206.07183v2}.
\end{thebibliography}
\end{document}
